\chapter{An audio file through the whole program}
\MARSChapterHeading{An audio file through the whole program}
\label{revision:wav}

The filter in Chapter~\ref{revision:filtering} accepts sample codes in memory.
A recording on disk contains more information: its sampling rate, the number
of channels, a description of the sample representation, and boundaries around
those bytes. To filter that recording, a caller must first establish where the
samples are and what they mean. It must then arrange them for the numerical
kernel and write a file that another program can interpret.

This chapter joins those operations in \texttt{wav-pipeline.asm}, with shared
functions in \texttt{wav-functions.asm}, both under \texttt{revision/labs/files}.
Keep the two files together: the caller includes the functions by filename.
The functions implement a deliberately restricted PCM WAV profile. We will
follow a four-frame file through parsing, filtering and writing, then use the
same program on a two-second generated input. The independent lab changes the
order of whole frames after filtering.

\section{Keep three counts distinct}

A sample code is one signed sixteen-bit amplitude. A \emph{frame} contains one
sample code for every channel at one instant. In stereo, the stored sequence is
left zero, right zero, left one, right one, and so on. At a sampling rate of
16000 frames per second, 32000 stereo frames describe two seconds, not four.
The PCM payload occupies $32000\cdot2\cdot2=128000$ bytes. A file also contains
headers and may contain unrelated chunks, so its byte count is larger.

Write $R$ for frames per second, $C$ for channel count and $N$ for frame count. The payload length is
$D=2CN$ bytes. Our accepted profile has $C\in\{1,2\}$, sampling rate 8000
through 48000, and $0\leq N\leq65536$. It accepts integer PCM format code one
with sixteen bits per sample. Compressed audio, floating-point samples and other
channel counts are outside this interface. Rejecting those inputs does not mean
that they are invalid WAV files in general.

\Needspace{15\baselineskip}
The parser produces the following little-endian descriptor. It is an internal
record, separate from the file's header layout.
\begin{table}[!ht]
  \centering
  \begin{tabular}{rll}
    \toprule
    Byte offset & Width & Meaning                                \\
    \midrule
    0           & 4     & Frames per second                      \\
    4           & 4     & Channel count $C$                      \\
    8           & 8     & Frame count $N$                        \\
    16          & 8     & PCM offset from the file buffer's base \\
    24          & 8     & PCM byte count $D$                     \\
    \bottomrule
  \end{tabular}
  \caption{The parser publishes all 32 descriptor bytes together.}
  \label{revision:wav-descriptor-layout}
\end{table}

The first sixteen bytes suffice to describe a separate PCM array. Accordingly,
\texttt{wav\_bound}, \texttt{wav\_filter} and \texttt{wav\_encode} consume only
that prefix. They do not follow the offset at byte sixteen. Once the caller has
computed a PCM pointer, adding the file offset again would select the wrong
input. The offset remains meaningful only relative to the original file buffer.

\section{Walk the file by its declared boundaries}

A RIFF file starts with four bytes spelling \texttt{RIFF}, a little-endian
unsigned word size, and four bytes spelling \texttt{WAVE}. The size excludes
the first eight bytes. In the trace input it is 74, giving 82 bytes in the whole
file. Our profile requires that exact extent: a shorter buffer is incomplete,
and extra bytes outside the declared image are rejected.

After these twelve bytes, each chunk has a four-byte identifier, an unsigned
four-byte payload length $L$, and $L$ payload bytes. An odd payload length has
one additional padding byte. The next chunk therefore starts at
\[
  q'=q+8+L+(L\bmod2),
\]
where $q$ is the current chunk's header offset. The padding belongs to the
chunk's storage extent but is excluded from its payload length.

\begin{table}[!ht]
  \centering
  \begin{tabular}{r l r r r}
    \toprule
    Header & Identifier                    & Payload start & Payload bytes & Next header \\
    \midrule
    12     & \texttt{JUNK}                 & 20            & 3             & 24          \\
    24     & \texttt{fmt\textvisiblespace} & 32            & 16            & 48          \\
    48     & \texttt{data}                 & 56            & 16            & 72          \\
    72     & \texttt{note}                 & 80            & 1             & 82          \\
    \bottomrule
  \end{tabular}
  \caption{Chunk boundaries in the 82-byte trace input. Offset 82 is the end,
    so there is no further header there.}
  \label{revision:wav-chunks}
\end{table}

The \texttt{fmt } identifier ends with an ASCII space. Its payload describes
PCM format, channels, rate, byte rate, block alignment and bits per sample.
Block alignment is the number of bytes in a frame, $2C$; byte rate must be
$2CR$. These redundant fields must agree. Our profile accepts a sixteen-byte
format payload or an eighteen-byte payload whose two-byte extension count is
zero. It requires exactly one format chunk before exactly one data chunk.
Duplicate format or data chunks are errors. Other chunks are skipped by extent,
including chunks after the data.

Searching memory for the four letters \texttt{data} would be incorrect. Those
letters could occur inside opaque payload bytes. Only an identifier reached by
the boundary walk identifies a chunk. Nor can the parser return as soon as it
finds a valid data chunk: the trailing \texttt{note} still needs its own valid
extent. In our trace, the odd padding bytes are \texttt{a5} and \texttt{7f}.
This local profile requires their presence without requiring their value to be
zero.

Check for eight available header bytes before reading the length. Decode the
unsigned word into a doubleword, then calculate the padded extent at that wider
width. A declared length of hexadecimal \texttt{ffffffff} needs 4294967296
bytes including padding. A word-sized addition would wrap that requirement to
zero. The parser uses the following sequence after decoding $L$ into
\texttt{\$t8}; \texttt{\$s3} holds $q$ and \texttt{\$s1} the complete file
byte count.
\Needspace{12\baselineskip}
\MARSCodeLabel{Include padding before comparing the next boundary}
\begin{minted}{mars-idea}
    andi $t7, $t8, 1
    daddu $t7, $t7, $t8
    daddiu $t4, $s3, 8
    daddu $t4, $t4, $t7
    sltu $t0, $s1, $t4
    bnezc $t0, parse_invalid
    nop
\end{minted}
The last comparison asks whether the required end lies beyond the buffer. It
is unsigned. At this stage the complete file count is already bounded by
266240 bytes, so these widened additions cannot overflow a doubleword.
Only a proven in-bounds payload can be interpreted or skipped.

After locating data, require $D$ to be divisible by $2C$ and derive $N=D/(2C)$.
The profile also limits non-PCM overhead to 4096 bytes, including headers and
padding. Thus the largest permitted complete file is
$4096+65536\cdot4=266240$ bytes. An empty recording is permitted: $N=D=0$,
with a valid format and an empty data chunk.

\section{Publish a descriptor only after the walk succeeds}

During parsing, keep candidate fields in registers or private storage. Finish
the boundary walk, verify all profile constraints, then copy the final record
to the caller's descriptor. If any test fails, all 32 original descriptor bytes
must remain unchanged. A plausible PCM region is not enough to publish a
partially trusted record.

For the trace input the result is rate 16000, two channels, four frames, offset
56 and length sixteen. In hexadecimal, 16000 is \texttt{3e80}; memory displays
its least significant byte first as \texttt{80 3e 00 00}. Likewise offset 56
appears as \texttt{38 00 00 00 00 00 00 00}. The native Debug view below shows
both halves of the completed record.
\begin{figure}[!ht]
  \centering
  \begin{tikzpicture}[alt={Two native Memory rows show rate 16000, channels two, frames four, PCM offset 56 and PCM length sixteen. A solid outline marks the sixteen-byte prefix and a dashed outline marks the offset and length.}]
    \node[anchor=north west,inner sep=0,font=\small\bfseries] (title) at (0,0) {A validated record, then its file-relative location};
    \node[anchor=north west,inner sep=0] (a) at ([yshift=-2mm]title.south west) {\MARSRevisionCrop{revision/captures/wav/wav-descriptor.png}{0.99\linewidth}{0.03125}{0.31555556}{0.58203125}{0.60888889}};
    \begin{scope}[shift={(a.south west)},x={($(a.south east)-(a.south west)$)},y={($(a.north west)-(a.south west)$)}]
      \draw[black,line width=.7pt] (.20,.3088) rectangle (.965,.6176);
      \draw[black,dashed,line width=.7pt] (.20,.005) rectangle (.965,.303);
    \end{scope}
  \end{tikzpicture}
  \caption{The solid outline contains the prefix consumed by the filter and writer. The dashed outline contains the file-relative offset and payload byte count. Each field uses its own width; neither row is a list of equal-sized words.}
  \label{fig:revision-wav-native-descriptor}
\end{figure}


Now change the trailing \texttt{note} length from one to three, keeping the
complete file at 82 bytes. Its payload would require bytes 80 through 82 and a
padding byte at 83. That exceeds the file even though format and PCM are intact.
The parser must return format error one and leave the caller's sentinel record
unchanged. This experiment distinguishes complete validation from returning
at the data chunk.
\begin{figure}[!ht]
\centering
\begin{tikzpicture}[alt={Two native Memory rows retain all thirty-two a5 sentinel bytes after the trailing chunk claims an extent beyond the complete input. The descriptor is unchanged.}]
\node[anchor=north west,inner sep=0,font=\small\bfseries] (title) at (0,0) {A late error leaves the record unpublished};
\node[anchor=north west,inner sep=0] (a) at ([yshift=-2mm]title.south west) {\MARSRevisionCrop{revision/captures/wav/wav-late-error.png}{0.99\linewidth}{0.03125}{0.31555556}{0.58203125}{0.60888889}};
\begin{scope}[shift={(a.south west)},x={($(a.south east)-(a.south west)$)},y={($(a.north west)-(a.south west)$)}]
\draw[black,line width=.7pt] (.20,.005) rectangle (.965,.6176);
\end{scope}
\end{tikzpicture}
\caption{All 32 descriptor bytes remain a5 after the malformed trailing chunk. The separately checked return status is one. Valid format and data chunks earlier in the file did not authorize publication.}
\label{fig:revision-wav-native-late-error}
\end{figure}


\Needspace{19\baselineskip}
\section{Give each channel its own time history}
The four input frames and their filtered results are:
\begin{table}[!ht]
  \centering
  \begin{tabular}{r rr rr}
    \toprule
          & \multicolumn{2}{c}{Input codes} & \multicolumn{2}{c}{Filtered codes}                  \\
    Frame & Left                            & Right                              & Left & Right   \\
    \midrule
    0     & 12000                           & $-12000$                           & 1500 & $-1500$ \\
    1     & 0                               & 16000                              & 1500 & 500     \\
    2     & $-8000$                         & 0                                  & 500  & 500     \\
    3     & 32767                           & $-32768$                           & 4596 & $-3596$ \\
    \bottomrule
  \end{tabular}
  \caption{An eight-tap average with seven initial zero samples in each channel.}
  \label{revision:wav-filter-trace}
\end{table}

Each of the eight coefficients is Q15 code 4096, representing one eighth.
The left output at frame two is $(12000+0-8000)/8=500$. At frame three,
$36767/8=4595.875$ rounds to 4596. The right calculation uses only right
samples: its frame-three sum is $-12000+16000+0-32768=-28768$, giving
$-3596$. No value clips.

Passing interleaved stereo bytes directly to a single-channel FIR would average
neighboring left and right samples. Instead, the wrapper gathers one channel
into a contiguous temporary window, calls the earlier kernel, and scatters
results back into the interleaved output. For frame $f$ and channel $c$, the
sample begins at byte displacement $2(fC+c)$. Here $0\leq c<C$.

The gather loop has block-start frame in \texttt{\$s6}, within-block index in
\texttt{\$t0}, channel count in \texttt{\$s2}, channel index in
\texttt{\$s5}, and original PCM base in \texttt{\$s0}:
\Needspace{10\baselineskip}
\MARSCodeLabel{Select one channel from an interleaved frame}
\begin{minted}{mars-idea}
    daddu $t1, $s6, $t0
    mul $t1, $t1, $s2
    daddu $t1, $t1, $s5
    dsll $t1, $t1, 1
    daddu $t1, $s0, $t1
\end{minted}
The word multiplication is safe for the validated frame and channel bounds:
$fC+c$ is at most 131071. The final addition to an address remains a
64-bit operation. Address width and index range are different arguments.

For each channel, initialize seven history codes to zero, then process blocks
of at most 64 new samples. A block of $m$ samples occupies a window of $7+m$
samples. After filtering, the seven values needed for the next call are the
last seven \emph{input} samples of that window. Their starting index is
$(7+m)-7=m$, a byte displacement of $2m$. This remains true for a short block:
if $m=3$, four of the retained values still come from the old history.
Retaining filtered outputs instead would introduce feedback and change the
specified FIR operation. Resetting history at each block would introduce a new
startup transient every 64 frames.

The kernel accumulates the exact eight products in doublewords, adds 16384,
shifts arithmetically right by fifteen, then clamps once. The wrapper totals
clipped sample codes across all channels and blocks. That count can reach
131072; it is neither a frame count nor a sixteen-bit quantity. Count zero
returns before reading PCM or coefficients, although the descriptor prefix
must still be valid.
\begin{figure}[!ht]
  \centering
  \begin{tikzpicture}[alt={One native Memory row holds four filtered stereo frames. Each frame contains a left and a right sixteen-bit sample. An outline marks the first frame, bytes dc 05 24 fa, or 1500 and minus 1500.}]
    \node[anchor=north west,inner sep=0,font=\small\bfseries] (title) at (0,0) {Frame boundaries remain after filtering};
    \node[anchor=north west,inner sep=0] (a) at ([yshift=-2mm]title.south west) {\MARSRevisionCrop{revision/captures/wav/wav-filter.png}{0.99\linewidth}{0.03125}{0.33888889}{0.58203125}{0.60888889}};
    \begin{scope}[shift={(a.south west)},x={($(a.south east)-(a.south west)$)},y={($(a.north west)-(a.south west)$)}]
      \draw[black,line width=.7pt] (.20,.01) rectangle (.39125,.467);
    \end{scope}
  \end{tikzpicture}
  \caption{The first outlined frame is left 1500, right minus 1500. The remaining frames are (1500,500), (500,500) and (4596,minus 3596). Each adjacent pair of bytes is one little-endian signed sample.}
  \label{fig:revision-wav-native-filter}
\end{figure}


\section{Write a canonical representation}

The output writer emits a twelve-byte RIFF/WAVE prefix, one sixteen-byte format
payload with its eight-byte chunk header, and a data header followed by the PCM.
The fixed overhead is $12+24+8=44$ bytes. Consequently
\texttt{wav\_bound} and a successful \texttt{wav\_encode} report
$B=44+2CN$. Our four-frame stereo output is sixty bytes; its RIFF size field
is 52. The input's unrelated chunks are deliberately omitted.

The file's channel count and block alignment are two-byte fields, while the
internal descriptor stores its channel count in four bytes. Do not copy the
internal prefix into a file header. Encode each required field at its specified
width and little-endian byte order. An empty recording still needs all 44
header bytes, with a data length of zero.

Validate the prefix and compare required $B$ with the unsigned output capacity
before writing any byte. Insufficient capacity returns status two and leaves
the destination unchanged; an invalid prefix returns one. A capacity with its
high bit set is large, rather than negative. The caller must still provide a
mapped destination interval of the actual required length. After successful
validation, copy exactly $D$ PCM bytes. No extra sample is needed for alignment.
\begin{figure}[!ht]
  \centering
  \begin{tikzpicture}[alt={Two native Memory rows begin at output offset thirty-two. The first twelve bytes finish the header; solid outlines mark sixteen PCM bytes spanning the rows. A dashed outline marks four following a5 guard bytes.}]
    \node[anchor=north west,inner sep=0,font=\small\bfseries] (title) at (0,0) {PCM starts after the canonical header};
    \node[anchor=north west,inner sep=0] (a) at ([yshift=-2mm]title.south west) {\MARSRevisionCrop{revision/captures/wav/wav-encoded.png}{0.99\linewidth}{0.03125}{0.31555556}{0.58203125}{0.60888889}};
    \begin{scope}[shift={(a.south west)},x={($(a.south east)-(a.south west)$)},y={($(a.north west)-(a.south west)$)}]
      \draw[black,line width=.7pt] (.77375,.3088) rectangle (.965,.6176);
      \draw[black,line width=.7pt] (.20,.005) rectangle (.77375,.303);
      \draw[black,dashed,line width=.7pt] (.77375,.005) rectangle (.965,.303);
    \end{scope}
  \end{tikzpicture}
  \caption{This view begins at output offset 32. The solid outlines start at offset 44 and contain exactly sixteen PCM bytes. The following four a5 bytes, dashed, lie outside the sixty-byte file and remain unchanged.}
  \label{fig:revision-wav-native-encoded}
\end{figure}


\section{Treat file handles and partial transfers as state}

\texttt{wav\_load(path,buffer,capacity,descriptor)} first checks capacity,
which must be between one and 266240. It then opens the input and accumulates
positive read counts. A short positive read is progress, not proof of EOF.
If the buffer becomes exactly full, a private one-byte probe distinguishes
an exact fit from an input that needs another byte. The probe does not write
past the caller's buffer. Only EOF establishes the complete input count;
parsing then decides whether that image is acceptable. Every acquired file
handle is closed, including on read or parse failure.

A failed load may have filled part of the byte buffer. Its guarantee concerns
the descriptor: no record is published for an incomplete or invalid input.
The caller must test the returned status before using either the record or any
assumed complete image. A failed load returns zero in \texttt{v1}.

\texttt{wav\_save(path,bytes,count)} validates the complete image before opening
an output that might truncate an existing file. Format rejection therefore
preserves that existing file. After opening, positive writes advance by exactly
the returned count until all selected bytes have been written. A zero-byte
write with data remaining is an error; retrying forever would make no progress.
A negative transfer result also fails. Close the acquired handle on either path.
An error after opening may leave a truncated or partially written output: this
interface does not implement a temporary-file rename transaction.

The numerical and file interfaces all return status in \texttt{v0}: zero for
success, one for format or profile error, two for a byte-count or capacity
error, and three for I/O failure. Successful \texttt{v1} has an interface-specific
unit: canonical bytes for bound, frames for parse, clipped sample codes for
filter, and complete file bytes for encode, load and save. Failed calls return
zero there. All preserve the full saved registers, frame pointer, stack pointer
and return continuation. Byte buffers and PCM may be unaligned; descriptor
records are eight-byte aligned. Selected input and output regions are disjoint
in the worked pipeline.

\section{Run a generated recording through the caller}

The fixture script creates deterministic bytes without replacing an existing
file. Start with the small trace so that the file boundaries and every output
sample remain inspectable. In the following commands, \texttt{CLI} denotes the
built MARS CLI JAR; use its full path when working elsewhere.
\Needspace{10\baselineskip}
\MARSCodeLabel{Create an input and prepare the pipeline controls}
\begin{Verbatim}
python3 packages/031-pcm-wav/fixtures/fixture.py input.wav --fixture trace
python3 packages/031-pcm-wav/prepare.py s63-1 \
  --source revision/labs/files/wav-pipeline.asm \
  --output revision/labs/files/wav-trace.asm
java -jar "$CLI" --start-at-main revision/labs/files/wav-trace.asm
\end{Verbatim}
The prepared source stays beside \texttt{wav-functions.asm}, so its include
continues to resolve. The paths in this recipe are \texttt{input.wav} and
\texttt{output.wav}, relative to the process's working directory. The preparer
changes caller data and controls while retaining the authored instructions.
It refuses to replace an existing output source. Use a fresh working directory
or move previous results before repeating an experiment.

Mode six loads, filters, encodes and saves in that order. Each stage tests
status before continuing. The global \texttt{clips} retains the filter's count;
\texttt{status} and \texttt{count} retain the final status and byte count.
For the trace, expect status zero, clips zero and count sixty. Break at
\texttt{filter\_return} to inspect PCM before the writer changes its
representation; break at \texttt{encode\_return} to inspect the complete
memory image before any file write. These caller boundaries have different
meanings even when the eventual output file is correct.

Next generate a fresh input with \texttt{--fixture listen} and prepare step
\texttt{s64-1}. This recording contains 32000 stereo frames at 16000 Hz.
Its input is 128056 bytes; its canonical filtered output is 128044 bytes.
The left channel contains a 200 Hz triangle pattern, the right a 250 Hz pattern,
and both include an alternating component at 8000 Hz with opposite signs.
An eight-sample average cancels that alternating component once its full
history is present. The first seven frames have the specified zero-history
transient. Listening may help connect codes to sound, but it cannot establish
rounding, channel order, clipping counts or memory bounds; compare those
numerically as well.

\section{Change frame order independently}

The lab's \texttt{wav-reverse-frames.asm} uses the same completed helpers and
adds one unfinished function. Reverse whole PCM frames after filtering and
before encoding. Each frame keeps its channels together, and each sixteen-bit
sample keeps its byte order. The trace output should begin with the filtered
frame $(4596,-3596)$ and finish with $(1500,-1500)$.

The function must also support exactly in-place reversal. Explain which original
bytes are needed before an endpoint store can overwrite them. An out-of-place
backward traversal and an in-place endpoint swap may use different loops.
Filtering after reversal would give a different signal because initial history
would attach to the other end; follow the caller's stated stage order.
The lab supplies the full independent contract and public behavioral checks.

\section*{Sources and laboratory connection}
The accepted profile and original six-function checks are in
\path{packages/031-pcm-wav/}. Microsoft and IBM's
\href{https://www.mmsp.ece.mcgill.ca/Documents/AudioFormats/WAVE/Docs/riffmci.pdf}{Multimedia
  Programming Interface and Data Specifications} describes the RIFF/WAVE
container and PCM fields. The profile restrictions, descriptor transaction and
kernel rounding policy in this chapter are local contracts. The written lab
uses the small trace to test those contracts before assigning frame reversal.
